\chapter{Análise de Variância (ANOVA)}
\label{cap:anova}

A análise de variância permite investigar diferenças entre médias de grupos a partir de um modelo para a resposta. Neste capítulo, vamos começar pela ANOVA de um fator, comparando as notas posteriores de Leitura, Vídeo e Oficina. Ela pode ser aplicada com dois ou mais níveis; é especialmente útil quando a pergunta envolve mais de dois grupos, pois não exige substituir a pergunta global por vários testes t separados \parencite{NISTAnova}.

\section{Da pergunta ao modelo}

No arquivo \texttt{exemplo\_notas.csv}, cada pessoa pertence a apenas uma atividade. A nota \texttt{depois} é a variável resposta contínua e \texttt{grupo} é o fator nominal. Para três grupos, a hipótese nula é:
\[
H_0:\mu_{\mathrm{Leitura}}=\mu_{\mathrm{Video}}=\mu_{\mathrm{Oficina}}.
\]
A alternativa afirma que pelo menos uma média difere de outra; ela não afirma que todos os pares sejam diferentes.

A lógica da ANOVA clássica compara a variação entre médias de grupos com a variação que permanece dentro dos grupos:
\[
F=\frac{QM_{\mathrm{entre}}}{QM_{\mathrm{dentro}}}.
\]
Os quadrados médios são somas de quadrados divididas pelos respectivos graus de liberdade. Assim, uma variação \textbf{entre grupos} grande em relação à variação \textbf{dentro dos grupos} produz um valor de $F$ maior. Grande dispersão dentro dos grupos, por si só, não favorece a detecção de diferenças entre médias.

\subsection{Condições que precisamos examinar}

A ANOVA clássica de grupos independentes usa um modelo com erros independentes, distribuição normal para a inferência exata e variância comum entre os grupos. Examine o desenho do estudo, os resíduos, os gráficos por grupo e possíveis valores influentes. A verificação não se resume a testar se a coluna inteira de notas é normal.

A ANOVA de Welch permite variâncias diferentes entre os grupos. Isso não corrige dependência entre observações nem torna a análise imune a valores extremos. Para medidas repetidas das mesmas pessoas, a estrutura do problema exige outro modelo; não empilhe notas anteriores e posteriores como se viessem de pessoas distintas.

\section{ANOVA de um fator no jamovi}

Desative os filtros usados no capítulo anterior e confira que o banco possui 18 pessoas, seis em cada grupo.

\begin{enumerate}
  \item Acesse \textbf{Análises $\rightarrow$ ANOVA $\rightarrow$ ANOVA de um fator (One-Way ANOVA)}.
  \item Mova \texttt{depois} para as variáveis dependentes e \texttt{grupo} para a variável de agrupamento.
  \item Solicite descritivas e o gráfico por grupo. Confira os valores da tabela~\ref{tab:descritivas-notas}.
  \item Escolha a versão clássica, identificada como \textbf{Fisher}, ou a versão \textbf{Welch}, de acordo com a análise planejada. O painel pode oferecer ambas; identifique qual resultado está relatando.
  \item Examine o gráfico Q-Q dos resíduos e a avaliação de homogeneidade, quando solicitados.
\end{enumerate}

As opções dessa análise, incluindo Welch e os procedimentos de comparação entre pares, constam na documentação oficial \parencite{JamoviAnova}.

\subsection{Conferindo o exemplo clássico}

Nos dados fictícios, a soma de quadrados entre grupos é aproximadamente $341{,}444$ e a soma de quadrados dentro dos grupos é $915$. Com três grupos e 18 observações, os graus de liberdade são $2$ e $15$. Portanto,
\[
F=\frac{341{,}444/2}{915/15}\approx2{,}799.
\]
O resultado clássico é $F(2,15)\approx2{,}799$, $p\approx0{,}093$. Com $\alpha=0{,}05$, não rejeitamos a igualdade das médias nesse teste. As médias amostrais são distintas, mas a informação disponível não permite declarar uma diferença populacional segundo a regra escolhida.

Para a versão de Welch, você deve encontrar aproximadamente $F=2{,}652$ e $p=0{,}119$. Os graus de liberdade do denominador são aproximados e diferem dos da ANOVA clássica. Nenhum dos resultados demonstra que as três atividades sejam equivalentes.

\section{Comparações múltiplas}

O teste global não identifica quais pares diferem. Quando comparações entre pares fizerem parte do plano de análise, precisamos considerar o problema da multiplicidade: realizar muitos testes sem ajuste aumenta a possibilidade de obter ao menos uma rejeição indevida.

No painel de comparações posteriores, o jamovi oferece \textbf{Tukey} para comparações com suposição de variâncias iguais e \textbf{Games--Howell} quando essa igualdade não é assumida \parencite{JamoviAnova}. Leia a diferença de médias e a direção da comparação, além do valor de $p$ ajustado. Um intervalo que inclua zero, quando produzido com o mesmo procedimento e nível, indica a incerteza correspondente àquela comparação.

No nosso exemplo, uma comparação Leitura menos Oficina por Tukey produz diferença de aproximadamente $-10{,}67$ pontos e $p$ ajustado de $0{,}077$. Essa comparação está apresentada apenas para praticar a leitura da saída: não selecione o par com o menor valor de $p$ e omita os demais. Contrastes definidos antes da coleta e explorações posteriores devem ser identificados no relatório.

\section{Magnitude e relato do resultado}

Para a ANOVA clássica de um fator, uma medida descritiva da magnitude é:
\[
\eta^2=\frac{SQ_{\mathrm{entre}}}{SQ_{\mathrm{total}}}.
\]
No exemplo, $SQ_{\mathrm{total}}=341{,}444+915$ e $\eta^2\approx0{,}272$. Essa é uma razão calculada sobre a amostra e o modelo; ela não indica que o método de ensino causou 27,2\% da variação das notas. Não transplante esse cálculo para a estatística de Welch como se fosse a mesma decomposição de somas de quadrados.

Um relato possível é: ``Nas 18 observações fictícias, as médias posteriores foram $59{,}17$, $64{,}67$ e $69{,}83$ pontos para Leitura, Vídeo e Oficina. A ANOVA clássica de um fator resultou em $F(2,15)=2{,}80$, $p=0{,}093$. O exemplo não forneceu evidência suficiente para rejeitar a igualdade das médias ao nível de 5\%.'' Em um trabalho real, acrescente os diagnósticos, as estimativas de interesse e a justificativa da análise escolhida.

\begin{tcolorbox}[colback=white,colframe=green,title={Dica de Conteúdo}]
ANOVA de grupos independentes, ANOVA de medidas repetidas e modelos com mais de um fator respondem a estruturas diferentes. Antes de adicionar fatores ou medidas, desenhe quem foi observado, quantas vezes e em quais condições. A organização do banco deve acompanhar esse desenho.
\end{tcolorbox}
